% This document was generated by an LLM.

\documentclass{essay}

\title{Continuity at a Point: Several Equivalent Languages and Their Subtleties}
\author{}
\date{}

\begin{document}

\maketitle

\section{Introduction}

Continuity is one of the central notions of real analysis, and it is also one of the first places where several apparently different mathematical languages turn out to describe exactly the same phenomenon. A function may be described as continuous at a point using inequalities involving $\varepsilon$ and $\delta$, using limits, using sequences, using neighborhoods, or, in a more abstract setting, using open sets and inverse images.

These formulations are closely related, but they emphasize different aspects of the same idea. The $\varepsilon$--$\delta$ definition expresses quantitative control. The limit formulation compares the behavior of the function near a point with its value at the point itself. The sequential formulation says that continuous functions preserve convergence. The neighborhood and topological formulations strip away coordinates and numerical distance and reveal continuity as a structural property.

The equivalence of these formulations can make continuity look simpler than it really is. Some details appear almost trivial when first encountered: whether the point $x=c$ is included in a condition, whether the domain contains points on both sides of $c$, whether a sequence must lie in the domain, or whether a point is isolated. These details are mathematically small but logically important. Ignoring them can lead to incorrect statements later.

This article compares the principal definitions of continuity at a point and pays particular attention to these subtle points.

\section{The Basic Setting}

Let
\[
f:D\to\mathbb{R},
\]
where $D\subseteq\mathbb{R}$, and let $c\in D$.

Informally, $f$ is continuous at $c$ if values of $f(x)$ can be made arbitrarily close to $f(c)$ by requiring $x$ to be sufficiently close to $c$. This statement has two ingredients:

\begin{itemize}
    \item closeness of inputs, measured by $|x-c|$;
    \item closeness of outputs, measured by $|f(x)-f(c)|$.
\end{itemize}

The purpose of a formal definition is to state precisely how these two kinds of closeness are related.

\section{The $\varepsilon$--$\delta$ Definition}

The standard direct definition is the following.

The function $f$ is continuous at $c$ if
\[
\forall \varepsilon>0\;\exists\delta>0\;\forall x\in D,
\qquad
|x-c|<\delta
\implies
|f(x)-f(c)|<\varepsilon.
\]

Equivalently,
\[
\boxed{
\forall \varepsilon>0\;\exists\delta>0\;\forall x\in D:
|x-c|<\delta
\implies
|f(x)-f(c)|<\varepsilon.
}
\]

The order of the quantifiers is essential. First an arbitrary output tolerance $\varepsilon>0$ is specified. Then one must find an input tolerance $\delta>0$ that works for every point $x$ of the domain satisfying $|x-c|<\delta$.

The number $\delta$ may depend on $\varepsilon$, but it may not depend on the particular choice of $x$ once $\varepsilon$ has been fixed. Thus continuity gives uniform control over all sufficiently nearby domain points.

The definition may be interpreted as a rule
\[
\text{desired output accuracy}
\longmapsto
\text{required input accuracy}.
\]

This is the most quantitative formulation of continuity.

\section{Continuity Expressed Through Limits}

A familiar alternative definition is
\[
\boxed{
\lim_{x\to c} f(x)=f(c).
}
\]

This formulation is especially natural in calculus. It says that the value approached by $f(x)$ as $x$ approaches $c$ is exactly the value that the function actually takes at $c$.

It is useful to separate these two pieces conceptually. The expression
\[
\lim_{x\to c} f(x)
\]
describes the behavior of $f$ near $c$, whereas
\[
f(c)
\]
describes the behavior of the function at the point $c$ itself.

Continuity asserts that these agree.

Thus one may summarize the idea as
\[
\boxed{
\text{continuous at }c
\iff
\text{limit near }c = \text{value at }c.
}
\]

Depending on the logical development of a textbook, this statement may be taken as the definition of continuity or proved as a theorem from the direct $\varepsilon$--$\delta$ definition.

\section{The Crucial Difference Between a Limit and Continuity}

The definition of a limit contains a subtle but important feature. The statement
\[
\lim_{x\to c} f(x)=L
\]
means that
\[
\forall \varepsilon>0\;\exists\delta>0
\]
such that
\[
0<|x-c|<\delta
\implies
|f(x)-L|<\varepsilon.
\]

The condition
\[
0<|x-c|
\]
means that the point $x=c$ is excluded.

This exclusion is not an accident. A limit describes what happens arbitrarily close to $c$, not what happens at $c$ itself. Therefore the value $f(c)$ is irrelevant to the existence or value of the limit.

For example, suppose
\[
f(x)=x^2 \qquad (x\neq 2),
\]
but define
\[
f(2)=1000.
\]
Then
\[
\lim_{x\to 2} f(x)=4,
\]
even though
\[
f(2)=1000.
\]
The function is therefore not continuous at $2$.

Now suppose we impose the condition
\[
\lim_{x\to c}f(x)=f(c).
\]
The $\varepsilon$--$\delta$ statement for the limit becomes
\[
0<|x-c|<\delta
\implies
|f(x)-f(c)|<\varepsilon.
\]

At the omitted point $x=c$, however,
\[
|f(c)-f(c)|=0<\varepsilon
\]
automatically. Therefore adding the point $x=c$ does not change the truth of the condition. We may replace
\[
0<|x-c|<\delta
\]
by
\[
|x-c|<\delta.
\]

This produces exactly the direct $\varepsilon$--$\delta$ definition of continuity.

Thus the tiny-looking distinction between
\[
0<|x-c|<\delta
\]
and
\[
|x-c|<\delta
\]
encodes a fundamental conceptual difference: limits ignore the value at the point, while continuity explicitly relates nearby behavior to the value at the point.

\section{The Sequential Definition}

For functions on the real numbers, continuity at a point can also be characterized entirely in terms of sequences.

The function $f$ is continuous at $c$ if and only if, for every sequence $(x_n)$ in the domain $D$,
\[
x_n\to c
\]
implies
\[
f(x_n)\to f(c).
\]

Symbolically,
\[
\boxed{
\forall (x_n)\subseteq D,
\qquad
x_n\to c
\implies
f(x_n)\to f(c).
}
\]

This says that continuous functions preserve convergence.

If
\[
x_n\to c,
\]
then applying a continuous function gives
\[
f(x_n)\to f(c).
\]

For example, because the function
\[
f(x)=x^2
\]
is continuous everywhere, from
\[
x_n\to 3
\]
we immediately obtain
\[
x_n^2\to 9.
\]

The sequential formulation is often easier to use than the $\varepsilon$--$\delta$ definition when studying convergence, because it allows one to reason with explicit objects rather than nested quantifiers.

\section{Sequences and Discontinuity}

The sequential criterion is particularly efficient for proving that a function is not continuous.

To show that $f$ is discontinuous at $c$, it is enough to find one sequence $(x_n)$ in the domain such that
\[
x_n\to c
\]
but
\[
f(x_n)\not\to f(c).
\]

For example, define
\[
f(x)=
\begin{cases}
0, & x<0,\\
1, & x\ge 0.
\end{cases}
\]

At $c=0$, consider the sequence
\[
x_n=-\frac{1}{n}.
\]
Then
\[
x_n\to 0,
\]
but
\[
f(x_n)=0
\]
for every $n$, while
\[
f(0)=1.
\]
Hence
\[
f(x_n)\not\to f(0),
\]
and $f$ is not continuous at $0$.

One must remember that the sequence must consist of points in the domain $D$. This becomes important when the domain is a proper subset of $\mathbb R$.

\section{The Neighborhood Definition}

The $\varepsilon$--$\delta$ definition may be reformulated geometrically using neighborhoods.

A neighborhood of a point $c$ is, informally, a set containing some open interval around $c$. Then $f$ is continuous at $c$ if for every neighborhood $V$ of $f(c)$, there exists a neighborhood $U$ of $c$ such that
\[
f(U\cap D)\subseteq V.
\]

Symbolically,
\[
\boxed{
\forall V\ni f(c)\;\exists U\ni c:
\quad
f(U\cap D)\subseteq V.
}
\]

For real-valued functions, one may take
\[
V=(f(c)-\varepsilon,f(c)+\varepsilon)
\]
and
\[
U=(c-\delta,c+\delta).
\]
Then the inclusion
\[
f(U\cap D)\subseteq V
\]
means exactly that
\[
|x-c|<\delta
\implies
|f(x)-f(c)|<\varepsilon.
\]

Thus the neighborhood formulation contains the same information as the $\varepsilon$--$\delta$ definition, but it no longer emphasizes numerical inequalities.

\section{The Inverse-Image Formulation}

The neighborhood definition can be rewritten in terms of inverse images.

For every neighborhood $V$ of $f(c)$, the inverse image
\[
f^{-1}(V)
\]
must contain a neighborhood of $c$ relative to the domain $D$.

One may write this as
\[
\boxed{
\forall V\ni f(c),
\qquad
c\in \operatorname{int}_D(f^{-1}(V)).
}
\]
where $\operatorname{int}_D$ denotes interior relative to the domain $D$.

The intuition is that every sufficiently small acceptable region around the output $f(c)$ pulls back to an acceptable region around the input $c$.

This inverse-image viewpoint becomes especially important in topology, where inverse images of open sets provide the standard global definition of continuity.

\section{Continuity in Metric Spaces}

The use of absolute values in the real definition is not essential. What is really needed is a notion of distance.

Let
\[
f:X\to Y
\]
where $X$ and $Y$ are metric spaces with metrics $d_X$ and $d_Y$. Then $f$ is continuous at $c\in X$ if
\[
\boxed{
\forall\varepsilon>0\;\exists\delta>0:
\quad
d_X(x,c)<\delta
\implies
d_Y(f(x),f(c))<\varepsilon.
}
\]

For the real numbers with the usual metric,
\[
d(x,y)=|x-y|,
\]
this becomes the familiar definition.

This shows that $\varepsilon$ and $\delta$ are not specifically tied to real numbers. They express a relation between distances in the domain and distances in the codomain.

\section{The Topological Formulation}

At an even more abstract level, distance itself can be discarded.

Let
\[
f:X\to Y
\]
be a function between topological spaces. Then $f$ is continuous at $c\in X$ if for every neighborhood $V$ of $f(c)$, the inverse image $f^{-1}(V)$ is a neighborhood of $c$.

Thus
\[
\boxed{
\forall V\ni f(c),
\qquad
f^{-1}(V)\text{ is a neighborhood of }c.
}
\]

This formulation contains no reference to numbers, absolute values, metrics, $\varepsilon$, or $\delta$. It expresses continuity purely in terms of the local structure of the spaces.

Continuity is therefore fundamentally a topological concept rather than merely an analytic one.

\section{Continuity of a Function on Its Entire Domain}

There is an important distinction between continuity at one point and continuity everywhere.

A function
\[
f:X\to Y
\]
is called continuous if it is continuous at every point $c\in X$.

For topological spaces, this has the elegant equivalent characterization
\[
\boxed{
f\text{ is continuous}
\iff
f^{-1}(U)\text{ is open in }X
\text{ for every open }U\subseteq Y.
}
\]

This is often taken as the definition of a continuous map in topology.

The formulation is global: instead of examining one point at a time, it describes how the function transforms all open sets by inverse image.

\section{Continuity Is Relative to the Domain}

One of the most easily overlooked subtleties is that continuity is always defined relative to the actual domain of the function.

Suppose
\[
f:[0,1]\to\mathbb R.
\]
To say that $f$ is continuous at $0$ means
\[
\forall\varepsilon>0\;\exists\delta>0:
\quad
x\in[0,1]
\text{ and }
|x|<\delta
\implies
|f(x)-f(0)|<\varepsilon.
\]

There is no requirement to consider negative $x$, because negative numbers are not in the domain.

Thus continuity at an endpoint of an interval is not defective or incomplete. The function is tested only against points that actually belong to its domain.

This observation also affects limits. If limits are defined relative to the domain $D$, then
\[
\lim_{x\to 0}f(x)
\]
for a function defined on $[0,1]$ is effectively a right-hand limit, because there are no domain points approaching $0$ from the left.

A common source of confusion is to silently replace the actual domain by all of $\mathbb R$. That changes the problem.

\section{Isolated Points}

The relative-domain viewpoint leads to an extreme but illuminating case.

Suppose $c$ is an isolated point of $D$. This means there exists some $\delta_0>0$ such that
\[
D\cap(c-\delta_0,c+\delta_0)=\{c\}.
\]

Then every function
\[
f:D\to\mathbb R
\]
is automatically continuous at $c$.

Indeed, let $\varepsilon>0$ be arbitrary and choose $\delta=\delta_0$. If $x\in D$ and
\[
|x-c|<\delta,
\]
then necessarily $x=c$. Therefore
\[
|f(x)-f(c)|=0<\varepsilon.
\]

No matter how strange the values of $f$ may be elsewhere, the function cannot fail to be continuous at an isolated point because there are no other domain points sufficiently close to that point.

For example, every point of $\mathbb N$ is isolated as a subset of $\mathbb R$. Hence every function
\[
f:\mathbb N\to\mathbb R
\]
is continuous at every point of its domain, regardless of the values it takes.

This example also exposes a subtlety in the slogan
\[
\lim_{x\to c}f(x)=f(c).
\]
If $c$ is not an accumulation point of the domain, then the punctured condition
\[
0<|x-c|<\delta
\]
may contain no domain points at all once $\delta$ is sufficiently small. Under a purely formal $\varepsilon$--$\delta$ reading, the limit condition would then be vacuously true for every proposed value of the limit. For this reason, many treatments define $\lim_{x\to c}f(x)$ only when $c$ is an accumulation point of the domain, or else handle isolated points as a separate convention.

This example emphasizes that continuity is a property of a function together with its domain, not merely of a symbolic formula.

\section{One-Sided Continuity}

For functions on subsets of $\mathbb R$, one may also consider one-sided versions of continuity.

The function $f$ is right-continuous at $c$ if
\[
\lim_{x\to c^+}f(x)=f(c),
\]
and left-continuous at $c$ if
\[
\lim_{x\to c^-}f(x)=f(c).
\]

If $c$ is an interior point of the domain and the domain contains points arbitrarily close to $c$ on both sides, then
\[
f\text{ is continuous at }c
\]
if and only if $f$ is both left-continuous and right-continuous at $c$.

At endpoints, however, only the side lying inside the domain is relevant to ordinary continuity relative to the domain.

\section{Sequential Continuity in General Topological Spaces}

For functions between metric spaces, the sequential characterization
\[
x_n\to c
\implies
f(x_n)\to f(c)
\]
is equivalent to continuity.

This remains true for the real numbers and for every standard space encountered in introductory analysis.

In completely arbitrary topological spaces, however, sequences are sometimes too weak to detect all of the topology. There exist topological spaces in which a function preserves the limits of every convergent sequence but is nevertheless not continuous.

In such settings, one needs more general objects such as nets or filters to recover a convergence-based characterization equivalent to topological continuity.

This distinction is not normally relevant in an introductory real-analysis course, but it is conceptually useful because it shows exactly how much structure the sequential criterion depends on.

\section{Comparison of the Main Formulations}

For a function
\[
f:D\subseteq\mathbb R\to\mathbb R
\]
and a point $c\in D$, the following conditions are equivalent.

First, the direct $\varepsilon$--$\delta$ condition:
\[
\forall\varepsilon>0\;\exists\delta>0:
\quad
|x-c|<\delta
\implies
|f(x)-f(c)|<\varepsilon.
\]

Second, the limit condition:
\[
\lim_{x\to c}f(x)=f(c),
\]
where the limit is interpreted relative to the domain.

Third, the sequential condition:
\[
x_n\to c
\implies
f(x_n)\to f(c)
\]
for every sequence $(x_n)$ contained in $D$.

Fourth, the neighborhood condition: every neighborhood of $f(c)$ contains the image of some sufficiently small neighborhood of $c$ relative to the domain.

Fifth, in topological language: every neighborhood of $f(c)$ has inverse image that is a neighborhood of $c$.

These formulations are not competing notions. They are different languages for the same local phenomenon.

\section{What Each Formulation Emphasizes}

Although equivalent in the usual setting of real analysis, the formulations are useful for different reasons.

The $\varepsilon$--$\delta$ definition emphasizes quantitative control:
\[
\text{output tolerance}
\longrightarrow
\text{input tolerance}.
\]

The limit formulation emphasizes agreement between nearby behavior and the actual function value:
\[
\text{behavior near }c
=
\text{behavior at }c.
\]

The sequential formulation emphasizes preservation of convergence:
\[
x_n\to c
\longrightarrow
f(x_n)\to f(c).
\]

The neighborhood formulation emphasizes local geometry:
\[
\text{near }c
\longrightarrow
\text{near }f(c).
\]

The topological formulation removes numerical distance altogether and reveals continuity as a structural property of spaces.

There is therefore a natural conceptual progression:
\[
\boxed{
\varepsilon\text{--}\delta
\longrightarrow
\text{metric spaces}
\longrightarrow
\text{neighborhoods}
\longrightarrow
\text{topology}.
}
\]

\section{Why the Small Logical Details Matter}

Continuity is a good example of a broader lesson in analysis: details that initially look pedantic often encode genuine mathematical distinctions.

The difference between
\[
0<|x-c|<\delta
\]
and
\[
|x-c|<\delta
\]
separates the notion of a limit from the notion of continuity.

The condition
\[
x\in D
\]
prevents us from testing the function at points where it is not defined.

The requirement that a sequence satisfy
\[
x_n\in D
\]
ensures that the sequential criterion respects the actual domain.

The concept of relative neighborhoods explains why endpoints cause no difficulty and why isolated points are automatically points of continuity.

Even the familiar identity
\[
\lim_{x\to c}f(x)=f(c)
\]
contains several logically distinct claims hidden behind compact notation: the limit must exist, the function value must be defined, and the two must agree.

These distinctions are easy to suppress once one becomes fluent with the notation, but keeping them visible while learning analysis prevents many later mistakes.

\section{Conclusion}

Continuity at a point can be expressed in several equivalent ways. In real analysis, the most important formulations are the $\varepsilon$--$\delta$ definition, the limit characterization, and the sequential criterion. Neighborhood language provides a bridge to metric spaces and topology, where continuity becomes a general structural notion independent of coordinates or numerical distance.

The equivalence of these definitions should not obscure their different conceptual roles. The $\varepsilon$--$\delta$ definition provides precise quantitative control. Limits separate nearby behavior from pointwise values. Sequences turn continuity into preservation of convergence. Neighborhoods and inverse images expose the local and topological structure behind the concept.

The small technical qualifications are part of the mathematics, not merely formal decoration. Whether $x=c$ is excluded, whether points lie in the domain, whether a point is isolated, and whether the surrounding space is metric all affect which statements are valid. Learning to notice such details is part of learning real analysis itself.

\end{document}
